\documentclass[11pt]{article}
\usepackage[a4paper,margin=1in]{geometry}
\usepackage[T1]{fontenc}
\usepackage[utf8]{inputenc}
\usepackage{lmodern}
\usepackage{amsmath,amssymb,mathtools}
\usepackage{microtype}
\setlength{\parindent}{1.2em}
\setlength{\parskip}{0.25em}
\newcommand{\dd}{\,d}
\newcommand{\modulo}[1]{\pmod{#1}}
\begin{document}

\begin{center}
{\Large G. Lejeune Dirichlet's Werke. II.}\par\vspace{1.2em}
{\Large XX.}\par\vspace{0.5em}
{\LARGE \textbf{Untersuchungen über ein Problem der Hydrodynamik}}\par\vspace{1em}
{\large Von G. Lejeune Dirichlet}\par\vspace{0.5em}
{\small Nachrichten von der G. A. Universität und der Königl. Gesellschaft der Wissenschaften zu Göttingen, Jahrg. 1857, Nr. 14, August 10, S. 205--207.}
\end{center}

\bigskip
\begin{center}
{\large \textbf{Untersuchungen über ein Problem der Hydrodynamik.}}
\end{center}

\begin{center}
[Auszug aus einer der Königlichen Societät am 31. Juli 1857 überreichten Abhandlung.\footnote{Der Auszug wird in Nr. 14 der Nachrichten der Göttinger Königlichen Gesellschaft der Wissenschaften auf Seite 205 mit den Worten eingeleitet: Am 31. Juli 1857 überreichte Hr. Prof. G. Lejeune Dirichlet der Königlichen Societät eine Abhandlung, die den Titel führt: ``Untersuchungen über ein Problem der Hydrodynamik'', von welcher das Nachfolgende ein Auszug ist. K.}]
\end{center}

Obgleich die Grundgleichungen der Hydrodynamik in der Form, in welcher man sie in den Lehrbüchern findet, seit Euler und in einer andern Form, welche man Lagrange verdankt, seit dem Erscheinen der ersten Ausgabe der \emph{Mécanique analytique} bekannt sind, so hat man doch in allen bis jetzt untersuchten Fällen, in welchen die Flüssigkeit im Laufe der Bewegung ihre Gestalt ändert, aus den Grundgleichungen nicht die vollständige Lösung des Problems abgeleitet, sondern sich auf eine genäherte Bestimmung der Bewegung beschränkt. Durch den eben erwähnten Umstand zu dem Versuche angeregt, ein Problem der bezeichneten Art in aller Strenge zu behandeln, war der Verfasser der Abhandlung bei der Schwierigkeit des Gegenstandes, dessen Erledigung die Integration eines Systems von nicht linearen partiellen Differentialgleichungen erfordert, darauf gefasst, dass ein solcher Versuch nur unter den einfachsten Voraussetzungen über die Bedingungen, welche die Bewegung der Flüssigkeit bestimmen, Erfolg haben würde, und deshalb nicht wenig überrascht, als sich nach einigen vergeblichen Bemühungen ein Fall darbot, der in so fern nicht zu den einfacheren zu zählen ist, als darin die gegenseitige Anziehung der Elemente der Flüssigkeit berücksichtigt wird, und in welchem ungeachtet dieses Umstandes die Bewegung nicht nur ihrem allgemeinen Charakter nach erkannt, sondern auch unter gewissen Beschränkungen in ihren Einzelheiten bestimmt werden kann.

Die Bedingungen dieses Falles der Bewegung, welcher in der Abhandlung behandelt wird, und das darauf bezügliche allgemeine Resultat sind in folgendem Satze ausgesprochen:

\begin{quote}
Hat eine homogene incompressible Flüssigkeit, die an ihrer Oberfläche einen constanten oder nur mit der Zeit veränderlichen Druck erleidet, anfänglich die Gestalt eines Ellipsoides; ist ferner die anfängliche Bewegung in zwei einfachere zerlegbar, eine Drehung der Flüssigkeit wie eines festen Körpers um eine durch den Schwerpunkt gehende Axe und eine zweite die relative Lage der Elemente ändernde Bewegung, bei welcher die Geschwindigkeiten derselben, senkrecht gegen drei bestimmte sich im Mittelpunkt rechtwinklig schneidende Ebenen zerlegt, den Abständen von den Ebenen proportional sind, so wird die Flüssigkeit in der in Folge eines solchen Anfangszustandes entstehenden Bewegung auch zu jeder späteren Zeit die Gestalt eines Ellipsoides haben, welches mit dem ursprünglichen concentrisch ist, dessen Axen sich aber im Allgemeinen mit der Zeit an Richtung und Grösse ändern. Von der augenblicklichen zu einer beliebigen Zeit stattfindenden Bewegung gilt dasselbe, was hinsichtlich der ursprünglichen vorausgesetzt worden ist, d. h. sie ist in zwei einfachere zerlegbar, wie sie vorhin definirt worden sind, so jedoch, dass sowohl die Drehungsaxe als die drei auf einander senkrechten Ebenen, auf welche sich diese Theilbewegungen beziehen, ebenfalls im Allgemeinen jeden Augenblick eine andere Lage annehmen.
\end{quote}

Zur vollständigen Kenntniss der Bewegung, deren allgemeine Beschaffenheit soeben angegeben worden ist, wird die Bestimmung von 9 Functionen der Zeit erfordert, welche durch eben so viel Gleichungen, eine endliche und 8 Differentialgleichungen zweiter Ordnung definirt werden. Von diesen letzteren lassen sich allgemein 7 Integrale erster Ordnung aufstellen, die jedoch zur vollständigen Lösung des Problems d. h. zur Zurückführung desselben auf Quadraturen nicht ausreichen, wenn nicht die anfängliche Gestalt und der im ersten Augenblick stattfindende Bewegungszustand, welcher 8 willkürliche Elemente einschliesst, weiteren Beschränkungen unterworfen werden.

Unter den auf Quadraturen zurückführbaren Fällen ist der einfachste, welcher allein hier erwähnt werden kann, der, wo die Masse ursprünglich die Form eines Umdrehungsellipsoides hat und keine anfänglichen Geschwindigkeiten vorhanden sind. Die Bewegung besteht dann aus isochronen Schwingungen, in welchen die Flüssigkeit durch die Kugelgestalt hindurchgehend abwechselnd die Form eines verlängerten und die eines abgeplatteten Ellipsoides annimmt.

\end{document}
